\section{Computability and the Refining Universe}

A discrete, local base is a \emph{computable} base. That is the deepest payoff of the framework's commitment to the eight atoms, but it is also the easiest claim to overstate. This section gives the careful version. It says exactly what computability rules out, exactly what it leaves open, and why the picture of a universe forever refining toward smoothness is a viewpoint rather than a destination.

Recall the base. The universe is a vibe mesh, a flow of eight atoms. The mesh is the whole $\{3,4,3,4\}$ weave. A dock is one cell, a here-now. Each dock carries exactly $24$ sites, the $24$ directions out of it, and each site holds one tone in $\{-1, 0, +1\}$, read as pain, peace, pleasure. The lean is the value order $-1 < 0 < +1$. The knit is the law of the state, collide then stream. The wake is the growing edge where new docks are born. The beat is the discrete step. Every one of these is finite and discrete. That single fact is what makes the universe a computation, and it is the fact this section weighs.

\subsection{What computability forbids, and what it does not}

Two plain definitions come first, because the whole argument turns on keeping them apart.

\begin{definition}[Computable]
A process is \emph{computable} when it could in principle be carried out, step by step, by an idealized computer. Each step is finite, each step is determined by a finite amount of local information, and the answer at any beat is reached after finitely many steps.
\end{definition}

\begin{definition}[Supertask]
A \emph{supertask} is the completion of infinitely many steps in a finite time, so as to read off a single answer that no finite prefix of the steps could give. No physical machine and no idealized computer can perform a supertask.
\end{definition}

With those in hand, the reach of computability is sharper than it first looks. It forbids two things and permits a third that is often confused with them.

It forbids the continuum. A genuine continuum is an actual smooth infinity of points packed between any two points, and it contains uncomputable real numbers, values that no program can name even in the limit. A computable base cannot host that. The sites are countable, the tones are three-valued, and there is no point between two adjacent docks.

It forbids supertasks. A computable base reaches each answer in finitely many beats. It never leans on the completed result of an infinite run. The knit advances the state one beat at a time, and the wake grows the space one beat at a time, and nothing in the flow waits on an infinite tally to finish.

It does \emph{not} forbid an infinite universe. An infinite mesh is allowed, so long as every finite region of it is finite and every step is local and finite. The number of docks may grow without bound, beat after beat, and the total may even be infinite in the limit, yet at no single beat is any computation infinite. Size is not the obstacle. Continuity and supertasks are.

\begin{principle}[Discrete-and-local, not finite-and-growing]
Computability does not push the theory toward a finite, ever-growing universe. It pushes toward a discrete, local one. Both the finite reading and the infinite reading pass the test, because both are discrete and both advance one finite beat at a time. The continuum and the supertask are what fail, and neither reading needs either.
\end{principle}

This is precisely why the two cosmic readings stay open and why the mechanisms by which infinity could emerge are allowed. Computability does not adjudicate between a bounded weave and an endless one. It only insists that whichever is true, it be made of countably many docks knit by a finite local law.

\subsection{The load-bearing commitment, discrete not continuous}

The base is discrete, top to bottom. The mesh is a weave of separate cells. The $24$ sites of a dock are a finite directed set. The tone is ternary. The beat is a tick, not a flow of instants. The symmetry of the cell is finite, the order-$1152$ group $F_4$, not a continuous Lie group. Nowhere in the eight atoms does a real-valued continuum appear.

Continuity, everywhere it shows up in the physics, is emergent. The continuous symmetry groups, the Dirac field, smooth Lorentz invariance, the inverse-square law of the attraction, all of these are coarse-grained limits of the discrete weave seen from far away. They are the look of a fine-enough fabric, never the threads themselves.

\begin{principle}[The discrete base]
The base is fully discrete. The mesh, the $24$ sites, the ternary tone, the discrete beat, and the finite symmetry are the whole of it. Continuity is always emergent, a coarse-grained limit and never a base ingredient. The universe is a computation, and the smoothness we observe is the appearance of a fine discrete weave viewed from a distance.
\end{principle}

The bridge between the two is emergent symmetry restoration. The finite cell symmetry sharpens, in the infrared, into the continuous group that the long-wavelength physics obeys. The discreteness never leaves. It merely stops being visible once the steps are small relative to the observer.

\subsection{The refining picture, and its honest limit}

It is tempting to dress the framework in a romantic image. The universe, one might say, refines toward the continuum as it grows, the weave forever tightening until it becomes smooth. The image is vivid and it is wrong, and the honest correction matters.

The large scale looks continuous because the discrete steps are tiny relative to us. It does not look continuous because the base is ever becoming continuous. A weave with ten docks across a meter and a weave with $10^{40}$ docks across a meter are both discrete. The second simply looks smoother to a being the size of a person. The gap between docks never closes to zero. It only shrinks below the scale at which we could notice it.

\begin{result}[The continuum is a viewpoint, not a destination]
The base stays discrete forever. The continuum is never reached and never approached as a limit of the substrate itself. It is a viewpoint, the appearance the weave takes on when the beat and the dock spacing fall far below the observer's resolution. Refinement, in the only honest sense, is the growth of scale separation, not a change in the kind of thing the base is.
\end{result}

So the refining universe is real as a statement about scale separation and false as a statement about ontology. The number of docks within any fixed coarse region can grow without bound, which is what sharpens the emergent continuous physics. The docks themselves stay countable and separate, forever.

\subsection{Why it matters}

Two consequences make the discreteness more than a tidy bookkeeping choice.

A computable cosmos is, in principle, fully simulable. There is no magic step, no uncomputable ingredient, no oracle hidden in the law. Given the eight atoms and the knit, an idealized computer could run the flow beat by beat. This is a strong honesty constraint on the theory. A model that needed a real-number continuum or a completed supertask to define even one of its moves would not be computable, and the framework deliberately needs neither.

And the discreteness frees the infinite-universe story. The worry about an infinite cosmos was never really about size. It was about the continuum smuggled in alongside it, the actual infinity of points, the uncomputable values, the supertasks lurking in any attempt to sum an endless smooth thing. Strip the continuum away and an infinite mesh is no harder to run, one finite beat at a time, than a finite one.

\begin{result}[Infinity was never the problem]
A computable cosmos is fully simulable, with no uncomputable ingredient. The same discreteness that secures simulability also frees the infinite reading, because infinity is permitted under computability while the continuum is not. The obstacle was always the continuum, never the size.
\end{result}

In the carrying image the fire is made of countably many sparks, however far it spreads. Each spark is a dock, each is finite, and the blaze is bright not because the sparks have merged into a smooth flame but because, from far enough away, countless separate sparks read as one continuous light.
